\documentclass[11pt,a4paper]{article}
\usepackage[T1]{fontenc}
\usepackage[utf8]{inputenc}
\usepackage{lmodern}
\usepackage[margin=21mm,top=19mm,bottom=20mm]{geometry}
\usepackage{amsmath,amssymb,booktabs,array,xcolor,microtype,url}
\definecolor{navy}{HTML}{17324D}
\definecolor{link}{HTML}{166F8A}
\definecolor{pale}{HTML}{EFF5F9}
\definecolor{rule}{HTML}{AEC4D3}
\setlength{\parindent}{0pt}
\setlength{\parskip}{6pt plus 1pt minus 1pt}
\setlength{\fboxsep}{8pt}
\setlength{\emergencystretch}{2em}
\clubpenalty=10000
\widowpenalty=10000
\pagestyle{plain}
\makeatletter
\renewcommand\section{\@startsection{section}{1}{0pt}{-2.0ex plus -.4ex minus -.2ex}{1ex plus .2ex}{\normalfont\Large\bfseries\color{navy}}}
\renewcommand\subsection{\@startsection{subsection}{2}{0pt}{-1.5ex plus -.3ex minus -.2ex}{.7ex plus .2ex}{\normalfont\normalsize\bfseries\color{navy}}}
\makeatother
\newcommand{\E}{\mathbb E}
\newcommand{\N}{\mathcal N}
\newcommand{\IG}{\operatorname{IG}}
\newcommand{\diag}{\operatorname{diag}}
\newcommand{\const}{\operatorname{const}}
\newcommand{\data}{\mathcal D}
\newcommand{\keybox}[1]{\par\smallskip\noindent\fcolorbox{rule}{pale}{\begin{minipage}{\dimexpr\linewidth-2\fboxsep-2\fboxrule\relax}#1\end{minipage}}\par\smallskip}
\newcommand{\href}[2]{\leavevmode\pdfstartlink attr{/Border[0 0 0]} user{/Subtype/Link/A<</S/URI/URI(\pdfescapestring{#1})>>}\textcolor{link}{#2}\pdfendlink}
\newcommand{\page}{\clearpage}
\begin{document}

{\small\color{link}\textbf{RESEARCH NOTE}\hfill Revised 30 October 2026}
\vspace{3pt}

{\LARGE\bfseries\color{navy} MCMC identification of a parametric Wiener-AR system\par}
{\large Arbitrary model orders, without sparsity-promoting priors\par}
\vspace{2pt}
\hrule height .5pt

\section{The model and the inference objective}

The aim is to infer the original linear-block, polynomial and output-noise
coefficients, together with the innovation variance. The orders
$n_a,n_b,n_f,n_d$ are arbitrary symbols throughout: no numerical values or
equality between them are assumed. For each identification they are specified
and held fixed; selecting their values is a separate model-selection problem.
We use proper Gaussian coefficient priors. This is a mathematical and
computational design for an estimator; no MCMC identification run is reported here.

With the delay operator $q^{-1}$, write
\begin{align}
 A(q^{-1};a)&=1+\sum_{i=1}^{n_a}a_iq^{-i}, &
 B(q^{-1};b)&=\sum_{i=1}^{n_b}b_iq^{-i}, \notag\\
 D(q^{-1};d)&=1+\sum_{i=1}^{n_d}d_iq^{-i}. \label{eq:polynomials}
\end{align}
The Wiener system and its autoregressive output noise are
\begin{align}
 \alpha_t&=\frac{B(q^{-1};b)}{A(q^{-1};a)}u_t,
 &\beta_t&=\alpha_t+\sum_{j=2}^{n_f}f_j\alpha_t^j, \notag\\
 c_t&=\beta_t+e_t,
 &D(q^{-1};d)e_t&=\nu_t,\quad \nu_t\overset{\mathrm{iid}}{\sim}\N(0,s),
 \qquad s=\sigma_\nu^2>0. \label{eq:model}
\end{align}
Inputs and innovations are assumed independent. There is no process noise inside
the linear signal block. Consequently, $\alpha$ and $\beta$ are deterministic
once the coefficients, input history and initial signal state have been specified.
Only the additive output-noise path is stochastic.

Keep the normalization $f_1=1$ and no constant polynomial term, as in the
manuscript. The physical coefficient vector is
\[
 \theta=(a_1,\ldots,a_{n_a},b_1,\ldots,b_{n_b},
              f_2,\ldots,f_{n_f},d_1,\ldots,d_{n_d})^\top.
\]
The normalization removes one scale ambiguity, assuming the true linear
polynomial coefficient is nonzero. Suitable excitation and a nonredundant model
are still needed; normalization alone does not prove identifiability.

\subsection{Dimensions and recursions for arbitrary orders}
There are $n_\theta=n_a+n_b+(n_f-1)+n_d$ free coefficients and one variance.
Thus the joint sampling dimension is $n_\theta+1$, before adding unknown
initial states. Here $n_f$ is the polynomial degree, while $n_f-1$ coefficients
are estimated because $f_1=1$. The direct recursions are
\begin{align}
 \alpha_t&=-\sum_{i=1}^{n_a}a_i\alpha_{t-i}
                 +\sum_{i=1}^{n_b}b_i u_{t-i}, \notag\\
 e_t&=-\sum_{i=1}^{n_d}d_i e_{t-i}+\nu_t. \label{eq:recursions}
\end{align}
Empty sums are zero: $n_a=0$ gives an FIR signal block and $n_d=0$ gives
white output noise. If $n_f=1$, there are no free nonlinear coefficients.

\keybox{For one identification, every chain and every posterior draw uses the
\textbf{same observed record} $\data=(u,c)$. MCMC iterations generate parameter
draws; they do not regenerate the input or the measured output. Physical time
$t$ and MCMC iteration $k$ are different indices.}

\page
\section{Specify initial conditions before writing the likelihood}

Initial conditions are part of the statistical model. They determine the first
samples of the predicted signal and the probability assigned to the beginning
of the noise record. Signal and noise initial conditions require separate choices.

\subsection{Initial conditions of the deterministic signal block}
For each proposed $(a,b)$, simulate the IIR block using one of the following
explicit models.
\begin{enumerate}
 \item \textbf{Known start and available input history.} Start from the known
 state before the experiment, propagate through all available input prehistory,
 and evaluate the likelihood only on the retained measured samples.
 \item \textbf{Unknown state at the beginning of the retained record.} Use a
 state-space realization of $B/A$, introduce its initial state $x_0$, and give
 it a proper Gaussian prior, for example $x_0\sim\N(0,V_x)$. Sample it together
 with $(\theta,s)$. The prior scale depends on the chosen state coordinates.
\end{enumerate}
A state-space initial state can absorb the effect of unobserved past inputs.
If a direct difference equation is used instead, the required past outputs and
past inputs must be supplied consistently.

The existing study generator runs a burn-in period and then discards those
samples. In general, the retained IIR trajectory begins with a nonzero state.
Resetting that state to zero defines a different forward model. Its transient
error may become small for long stable records, but that is an approximation,
not an exact likelihood for the generated record.

\subsection{Initial conditions of the AR output noise}
Three common choices are legitimate, but they must not be mixed:
\begin{itemize}
 \item \textbf{Stationary start:} the initial noise state follows the stationary
 distribution determined by $(d,s)$. The main worked likelihood in this note
 uses this choice.
 \item \textbf{Known zero or known nonzero start:} condition on the specified
 noise state and propagate the innovations from that start.
 \item \textbf{Conditional likelihood:} condition on an initial block of
 observations and omit its density. This uses fewer likelihood contributions.
\end{itemize}
A finite burn-in from zero can approximate a stationary start. It is not
identical to stationarity, particularly for poles close to the unit circle.
For an exact simulation benchmark, match the likelihood to the actual generator.

\subsection{General Gaussian form}
Let $e=c-\beta(\theta)$, suppressing a known or sampled $x_0$. For a zero-mean
noise initialization,
\begin{equation}
 c\mid u,\theta,s\sim\N\!\left(\beta(\theta),sK(d)\right),\qquad
 Q=e^\top K(d)^{-1}e. \label{eq:matrixlikelihood}
\end{equation}
Hence
\begin{equation}
 \log L=-\frac N2\log(2\pi s)-\frac12\log|K(d)|-\frac{Q}{2s}.
 \label{eq:generalloglikelihood}
\end{equation}
For known zero noise prehistory, $K=HH^\top$ with a unit-diagonal lower triangular
convolution matrix $H$, so $|K|=1$; this finite-record covariance is generally
not Toeplitz. The stationary covariance is Toeplitz and includes uncertainty
in the initial noise state. Evaluate quadratic forms by solves, never by
explicitly forming $K^{-1}$. For a known nonzero initial noise state, add its
propagated mean to $\beta$ before forming the residual vector.

\page
\section{How the function $Q(\theta)$ is formed}
\label{sec:Qconstruction}

$Q$ measures the discrepancy between the observed record and a proposed
noise-free response, with the residual correlation accounted for. It is rebuilt
at every proposed $\theta$; it is not a fixed function of the observations alone.
For clarity, the following formulas suppress a known or sampled signal initial
state $x_0$. If it is unknown, write $Q(\theta,x_0)$ instead.
$Q$ is quadratic in the residual vector for fixed $d$, but is generally not
quadratic in the full parameter vector $\theta$.

\subsection{Step 1: produce the proposed noise-free response}
Use $(a,b)$ and the fixed input to simulate the first recursion in
\eqref{eq:recursions}, then apply the polynomial:
\[
 \beta_t(\theta)=\alpha_t(a,b)+\sum_{j=2}^{n_f}f_j\alpha_t(a,b)^j.
\]
Changing $a$ or $b$ changes the full internal trajectory, not only one
regression coefficient. Changing $f$ changes the static transformation of it.

\subsection{Step 2: subtract it from the observations}
Form the output residual vector
\[
 e_t(\theta)=c_t-\beta_t(\theta),\qquad
 e(\theta)=(e_1(\theta),\ldots,e_N(\theta))^\top.
\]
The unweighted sum $\sum_t e_t^2$ ignores serial correlation. For AR output
noise, neighboring residuals must be interpreted jointly through $d$.

\subsection{Step 3: whiten the residuals using the proposed noise model}
At times with all required lags inside the record, apply $D(q^{-1};d)$:
\begin{equation}
 \begin{aligned}
 r_t(\theta)&=e_t(\theta)+\sum_{i=1}^{n_d}d_i e_{t-i}(\theta)\\
 &=c_t-\beta_t(\theta)+\sum_{i=1}^{n_d}d_i
                  [c_{t-i}-\beta_{t-i}(\theta)].
 \end{aligned} \label{eq:innovationresidual}
\end{equation}
These are candidate innovation residuals. At the true coefficients and with
the correct boundary treatment, they recover the driving innovations. Away
from those coefficients they generally contain both model error and noise.

\subsection{Step 4: sum squared whitened residuals, including the boundary}
In the matrix formulation, factor $K(d)=L_dL_d^\top$ by Cholesky and solve
$L_dw=e(\theta)$. Then
\begin{equation}
 \boxed{Q(\theta)=e(\theta)^\top K(d)^{-1}e(\theta)
                     =\|L_d^{-1}e(\theta)\|_2^2=\sum_{t=1}^N w_t^2.}
 \label{eq:Qmatrix}
\end{equation}
The next section gives the same quantity using only a small initial block and
the scalar residuals $r_t$, avoiding a dense $N\times N$ factorization.

\keybox{$K(d)$ is the covariance for \textbf{unit innovation variance}.
Consequently $Q(\theta)$ does not contain $s$: the full Gaussian exponent uses
$Q(\theta)/(2s)$. Both the residuals (through $a,b,f$) and their whitening
(through $d$) change with $\theta$. For white noise, $K=I$ and $Q=\sum_t e_t^2$.}

\page
\section{Compute $Q$ and the likelihood for arbitrary AR order}
\label{sec:innovlikelihood}

For stationary output noise, define $\gamma_k(d)=\operatorname{Cov}(e_t,e_{t-k})/s$.
These unit-variance autocovariances depend on $d$, not on $s$. For $n_d>0$,
the Yule-Walker equations are
\begin{equation}
 \gamma_0+\sum_{i=1}^{n_d}d_i\gamma_i=1,\qquad
 \gamma_k+\sum_{i=1}^{n_d}d_i\gamma_{|k-i|}=0,
 \quad k=1,\ldots,n_d. \label{eq:yulewalker}
\end{equation}
Solve for $\gamma_0,\ldots,\gamma_{n_d}$. Let $n_0=\min(n_d,N)$ and form the
initial-block covariance
\begin{equation}
 P_d=[\gamma_{|i-j|}(d)]_{i,j=1}^{n_0},\qquad
 e_{1:n_0}\sim\N(0,sP_d). \label{eq:Pd}
\end{equation}
A small state-space Lyapunov equation is an equivalent way to compute these
covariances. They must be recomputed at each proposed $d$.

\subsection{Stationary initial block plus subsequent innovations}
Factor the residual density into its initial block and any subsequent AR
transition densities. With $P_d=L_PL_P^\top$, this yields
\begin{equation}
 \boxed{\begin{aligned}
 Q(\theta)&=\underbrace{\|L_P^{-1}e_{1:n_0}(\theta)\|_2^2}_{\text{initial-block contribution}}
             +\underbrace{\sum_{t=n_0+1}^N r_t(\theta)^2}_{\text{innovation contribution}},\\
 \log L(\theta,s)&=-\frac N2\log(2\pi s)-\frac12\log|P_d|
                                  -\frac{Q(\theta)}{2s}.
 \end{aligned}} \label{eq:Qstationary}
\end{equation}
The first term assigns the correct probability to the beginning of the record,
whose missing earlier noise samples are integrated out. Simply setting them
to zero would give a different $Q$. The transformation from $e$ to
$(e_{1:n_0},r_{n_0+1:N})$ is triangular with unit diagonal. Thus
$|K(d)|=|P_d|$, and \eqref{eq:Qstationary} equals the matrix Gaussian likelihood.

Compute the first contribution by solving $L_Pw_0=e_{1:n_0}$ and taking
$w_0^\top w_0$. Then accumulate $r_t^2$ sequentially. Also compute
$\log|P_d|=2\sum_i\log (L_P)_{ii}$. A fixed small AR order gives a cost
linear in $N$, plus the small covariance calculation.

\subsection{Other boundary conventions}
With \textbf{known zero noise prehistory}, set $e_t=0$ for $t\leq0$, form
$r_t=e_t+\sum_i d_i e_{t-i}$ for every $t=1,\ldots,N$, and use
$Q=\sum_{t=1}^N r_t^2$. The log determinant is zero. With a \textbf{conditional}
likelihood given the first $n_0$ outputs, omit the initial-block contribution,
use $Q_{\mathrm{cond}}=\sum_{t=n_0+1}^N r_t^2$, and replace $N$ by $N-n_0$
in the Gaussian normalization. Neither choice is the full stationary likelihood.

For $n_d=0$, the initial block is empty, its contribution and log determinant
are zero, and $r_t=e_t$. If $N\leq n_d$, the initial block contains the whole
record and the subsequent-innovation sum is empty.

\keybox{$Q$ is the squared-error part of the likelihood, not the whole
likelihood. When $d$ and $s$ are unknown, retain both the covariance determinant
and the $-\tfrac N2\log s$ term. Differentiation must include the dependence of
$P_d$ on $d$ and of $e$ on the signal-path coefficients.}

\page
\section{Reasonable priors without sparsity promotion}

Choose proper, fixed covariance matrices, in the physical units of the
coefficients. A starting family is
\begin{align}
 p(a)&\propto\N(a;0,V_a)\,\mathbf1\{A\text{ is stable}\}, &
 b&\sim\N(0,V_b), \notag\\
 f&\sim\N(0,V_f), &
 p(d)&\propto\N(d;0,V_d)\,\mathbf1\{D\text{ is stable}\}, \notag\\
 s&\sim\IG(\xi_0,\zeta_0), &&\xi_0,\zeta_0>0. \label{eq:priors}
\end{align}
Here $f=(f_2,\ldots,f_{n_f})^\top$. These parameter blocks are independent
\emph{a priori}, including independence from $s$. Within a stability-truncated
block, the coefficients need not be independent even if the underlying
Gaussian covariance is diagonal. Nonzero prior means can be used when justified.
Empty parameter blocks are omitted, with zero contribution to the log prior.

The variance convention is
\begin{equation}
 p(s)=\frac{\zeta_0^{\xi_0}}{\Gamma(\xi_0)}
           s^{-\xi_0-1}\exp(-\zeta_0/s),\quad s>0.
 \label{eq:ig}
\end{equation}
Equivalently, the precision $\tau=1/s$ has a Gamma prior with shape $\xi_0$
and \textbf{rate} $\zeta_0$. A Gamma prior directly on $s$ is also valid, but
does not give the usual conjugate Gaussian variance update.

\subsection{An illustrative calibration, not a universal default}
For arbitrary block lengths, possible covariance families are
\[
 V_a=\tau_a^2I_{n_a},\quad V_b=\tau_b^2I_{n_b},\quad
 V_f=\diag(\tau_{f,2}^2,\ldots,\tau_{f,n_f}^2),\quad
 V_d=\tau_d^2I_{n_d},
\]
with fixed positive scales chosen in the coefficient units. The Gaussian
densities for $a$ and $d$ remain restricted to stable filters; individual
coefficient bounds alone do not enforce higher-order stability. Distinct
polynomial scales can reflect different plausible contributions of different
powers over the expected input range. For the variance, one possible choice is
\[
 s\sim\IG(3,2\sigma_{\mathrm{ref}}^2),\qquad
 \E[s]=\sigma_{\mathrm{ref}}^2,\quad
 \operatorname{SD}(s)=\sigma_{\mathrm{ref}}^2.
\]
Choose $\sigma_{\mathrm{ref}}$ from independent calibration, measurement
specifications or a defensible noise range. Do not use the hidden true variance
of a simulation as information available to the estimator. These scales should
be revised if their implied outputs are implausible.

Draw coefficients and variances from the prior, simulate responses for the
experimental input, and inspect amplitudes, transients and noise correlation.
This prior predictive check is more informative than inspecting coefficient
standard deviations in isolation. Stable poles can still imply very large gains.

\subsection{Units and interpretation}
Fixed Gaussian priors regularize the fit but do not select a sparse model:
no coefficient is pruned and no coefficient-specific precision is learned.
An excessively tight prior can materially affect finite-sample estimates.

If rescaling is useful, preserve the normalization. Under $u'=u/s_u$,
$c'=c/s_c$ and $\alpha'=\alpha/s_c$,
\[
 a'=a,\quad d'=d,\quad b'=\frac{s_u}{s_c}b,\quad
 f'_j=f_js_c^{j-1},\quad s'=s/s_c^2.
\]
Transform priors consistently. Arbitrarily subtracting the output mean is not
innocuous when the static model has no constant term.

\page
\section{Enforce stability with smooth sampling coordinates}
\label{sec:coordinates}

This section separates three questions: which physical parameters are allowed,
how an unrestricted numerical vector represents them, and how their probability
density must change when the coordinates change. The model orders remain
arbitrary throughout.

\subsection{Physical parameters and numerical coordinates are different objects}
The physical unknowns are $x=(\theta,s)$: the coefficients in the Wiener-AR
equations and the innovation variance. The prior already restricts $a$ and $d$
to stable denominator polynomials and requires $s>0$. We now seek a smooth,
one-to-one map $x=g(z)$ from unrestricted real coordinates $z$ onto that domain.
The transformation implements the chosen domain; it does not add parameters
or change the forward equations.

For a generic denominator with $m$ free coefficients, write
\[
 H(q^{-1};h)=1+\sum_{j=1}^m h_jq^{-j},\qquad
 C_m(\lambda)=\lambda^m+h_1\lambda^{m-1}+\cdots+h_m.
\]
Stability means that the roots of $C_m$ lie strictly inside the unit circle.
The coefficients $h_j$ are coupled by this condition. Restricting each
$h_j$ separately to $(-1,1)$ does not, in general, describe the stable region.
We use $m=n_a$ and $h=a$ for the signal denominator, then $m=n_d$ and $h=d$
for the noise denominator. Neither order is assigned a numerical value.

\subsection{First map: real coordinates to reflection coefficients}
Introduce an auxiliary vector $\kappa=(\kappa_1,\ldots,\kappa_m)^\top$.
These reflection coefficients are coordinates for a stable polynomial;
they are generally different from its physical coefficients $h$.
For every real $z_{c,i}$, set
\begin{equation}
 \kappa_i=\tanh z_{c,i}
   =\frac{\exp(2z_{c,i})-1}{\exp(2z_{c,i})+1}\in(-1,1),\qquad
 \frac{\partial\kappa_i}{\partial z_{c,i}}=1-\kappa_i^2>0.
 \label{eq:tanhmap}
\end{equation}
The inverse is
\[
 z_{c,i}=\operatorname{atanh}\kappa_i
      =\frac12\log\frac{1+\kappa_i}{1-\kappa_i}.
\]
Thus every interior value $\kappa_i\in(-1,1)$ corresponds to exactly one
finite real coordinate. No clipping or boundary rejection is needed to impose
these interval constraints mathematically. The next map turns $\kappa$ into $h$.

The positive innovation variance uses a separate scalar transformation:
\begin{equation}
 s=\exp v,\qquad v=\log s\in\mathbb R,\qquad \frac{ds}{dv}=\exp v=s.
 \label{eq:variancemap}
\end{equation}
The numerator coefficients $b$ and polynomial coefficients $f$ stay in their
original coordinates; their transformation matrices are identities.

\keybox{The numerical sampler will move in $z$. At every proposed $z$, we
recover the physical $(a,b,f,d,s)$ and evaluate the same Wiener-AR model and
physical priors. The maps $\tanh$ and $\exp$ are computational choices; the
prior on the physical parameters remains the one specified in Section 5.}

\page
\subsection{Second map: reflection coefficients to the physical polynomial}

The Schur/Levinson step-up recursion builds the denominator one order at a time
[7]. Here $r$ is the intermediate polynomial order, not a data-time index or
an MCMC iteration. The superscript in $h_j^{(r)}$ denotes that order, not a power.

Start with an empty coefficient vector $h^{(0)}$. For each $r=1,\ldots,m$,
take a copy of all coefficients at order $r-1$ and compute
\begin{equation}
 \begin{aligned}
 h_r^{(r)}&=\kappa_r,\\
 h_j^{(r)}&=h_j^{(r-1)}+\kappa_rh_{r-j}^{(r-1)},
                   \qquad j=1,\ldots,r-1.
 \end{aligned} \label{eq:schur}
\end{equation}
The new coefficient is $\kappa_r$. Each earlier coefficient receives a correction
equal to $\kappa_r$ times the coefficient at its reversed position in the
\emph{previous} vector. All right-hand sides use that previous vector, so an
implementation must not overwrite it partway through the step.

After the final step, $h=h^{(m)}$; denote this result by $\mathcal T_m(\kappa)$.
The real Schur stability criterion states that the resulting polynomial is
strictly stable exactly when all its reflection coefficients satisfy
$|\kappa_i|<1$ [7]. This is the reason for bounding $\kappa$, rather than
independently bounding the original polynomial coefficients.

\subsection{Why the map can be inverted}
At a generic recursion step, a reversed pair obeys
\[
 h_j^{(r)}=u_j+\kappa_ru_{r-j},\qquad
 h_{r-j}^{(r)}=u_{r-j}+\kappa_ru_j,
 \qquad u=h^{(r-1)}.
\]
Subtract $\kappa_r$ times the second equality from the first. The terms in
$u_{r-j}$ cancel, leaving $(1-\kappa_r^2)u_j$. Hence the step-down inverse is
\begin{equation}
 \kappa_r=h_r^{(r)},\qquad
 h_j^{(r-1)}=
 \frac{h_j^{(r)}-\kappa_rh_{r-j}^{(r)}}{1-\kappa_r^2},
 \quad j=1,\ldots,r-1. \label{eq:schurinverse}
\end{equation}
Since $|\kappa_r|<1$, the denominator is positive. Repeating backwards recovers
all reflection coefficients, then $\operatorname{atanh}$ recovers $z_c$.
This explains the one-to-one correspondence in the interior of the stable domain.

\subsection{Assemble the complete coordinate map}
Apply the same generic construction independently to $a$ and $d$:
\begin{equation}
 a=\mathcal T_{n_a}(\tanh z_a),\qquad
 d=\mathcal T_{n_d}(\tanh z_d),\qquad s=\exp v.
 \label{eq:transform}
\end{equation}
The sampling vector is
\[
 z=(z_a^\top,b^\top,f^\top,z_d^\top,v)^\top\in\mathbb R^{n_\theta+1}.
\]
Unknown signal initial states can be appended in their original real
coordinates. If a denominator has order zero, its coordinate block and recursion
are empty. The number of unknowns is unchanged by this transformation.

\page
\subsection{Why a change of coordinates requires a density correction}

Let $\pi_x(x)=p(\theta,s\mid\data)$ denote the intended posterior density with
respect to physical-coordinate volume $dx$. Consider a small region around
$z$. Its image under $x=g(z)$ has volume
\[
 dx\approx\left|\det Dg(z)\right|\,dz,\qquad
 Dg(z)=\frac{\partial x}{\partial z}.
\]
The same event must have the same probability in either set of coordinates:
\begin{align}
 \pi_z(z)\,dz&\approx\pi_x(g(z))\,dx \notag\\
             &\approx\pi_x(g(z))\,|\det Dg(z)|\,dz. \notag
\end{align}
Dividing by $dz$ and taking the limit gives the change-of-variables formula [2]:
\begin{equation}
 \boxed{\pi_z(z)=\pi_x(g(z))\,J(z),\qquad J(z)=|\det Dg(z)|.}
 \label{eq:densitychange}
\end{equation}
The factor is multiplied because we used the derivative of the map
\emph{from sampling coordinates to physical coordinates}. Using the inverse
derivative would give the reciprocal factor. In log form,
\[
 \log\pi_z(z)=\log\pi_x(g(z))+\log J(z).
\]
The probability distribution is preserved, although its density height changes
with the local volume scale. If one samples $\pi_x(g(z))$ without the factor,
the mapped physical draws instead have density proportional to
$\pi_x(x)/J(g^{-1}(x))$. The missing factor therefore changes the answer.

\subsection{A scalar illustration of the direction of the correction}
Consider one auxiliary scalar $\kappa$ with a uniform density $1/2$ on $(-1,1)$.
This is an illustration of coordinate volume, not a proposed prior for the
Wiener coefficients. If $\kappa=\tanh z_c$, then
\[
 p_{z_c}(z_c)=\frac12(1-\tanh^2 z_c).
\]
Indeed, its integral over $\mathbb R$ is
$\tfrac12[\tanh z_c]_{-\infty}^{\infty}=1$. Equal-length intervals far from
zero map to shorter intervals near $\kappa=\pm1$, so they must receive less
probability. Omitting the derivative would leave a constant over $\mathbb R$,
which cannot even be normalized.

\subsection{A worked correction for the innovation-variance prior}
The prior is $p_s(s)\propto s^{-\xi_0-1}\exp(-\zeta_0/s)$. Substitute $s=e^v$
and multiply by $ds/dv=e^v$:
\begin{equation}
 \begin{aligned}
 p_v(v)&=p_s(e^v)e^v
       \propto e^{-(\xi_0+1)v}\exp(-\zeta_0e^{-v})e^v,\\
 \log p_v(v)&=\const-\xi_0v-\zeta_0e^{-v}.
 \end{aligned} \label{eq:variancepriorchange}
\end{equation}
Thus the extra $+v$ removes one power of $s$ from the density expression.
The physical prior remains inverse-gamma; its density in the log-variance
coordinate has a different functional form.

\page
\subsection{Derive the Jacobian of one denominator map, step by step}

The denominator map is the composition $z_c\mapsto\kappa\mapsto h$.
By the chain rule for matrices and multiplicativity of determinants,
\begin{equation}
 \mathcal J_m(z_c)
 =\left|\det\frac{\partial h}{\partial\kappa}\right|
  \left|\det\frac{\partial\kappa}{\partial z_c}\right|.
 \label{eq:jacchain}
\end{equation}
\textbf{The first coordinate map.} Each $\kappa_i$ depends only on $z_{c,i}$,
so \eqref{eq:tanhmap} gives a diagonal matrix and its determinant:
\[
 \frac{\partial\kappa}{\partial z_c}
      =\diag(1-\kappa_1^2,\ldots,1-\kappa_m^2),\qquad
 \det\frac{\partial\kappa}{\partial z_c}
      =\prod_{i=1}^m(1-\kappa_i^2).
\]
\textbf{One recursion step.} Put $u=h^{(r-1)}$ and let $R_{r-1}$ reverse the
order of its entries. Equation~\eqref{eq:schur} can be written
\[
 h^{(r)}=
 \begin{pmatrix}(I_{r-1}+\kappa_rR_{r-1})u\\\kappa_r\end{pmatrix}.
\]
Differentiate with respect to the input vector $(u^\top,\kappa_r)^\top$:
\begin{equation}
 \frac{\partial h^{(r)}}{\partial(u,\kappa_r)}=
 \begin{pmatrix}
 I_{r-1}+\kappa_rR_{r-1}&R_{r-1}u\\
 0^\top&1
 \end{pmatrix}. \label{eq:recursionjacobian}
\end{equation}
It is block triangular. Its determinant is therefore
$\det(I_{r-1}+\kappa_rR_{r-1})$; the column $R_{r-1}u$ does not affect it.

\textbf{Evaluate that determinant.} Reversing twice returns the original
vector, so $R_{r-1}^2=I$. Its eigenvectors can be symmetric or antisymmetric
under reversal. Each mirrored pair of coordinates gives one of each kind; if
there is an unpaired middle coordinate, it is symmetric. Consequently there
are $\lceil(r-1)/2\rceil$ eigenvalues $+1$ and $\lfloor(r-1)/2\rfloor$
eigenvalues $-1$. The eigenvalues of $I+\kappa_rR$ are thus $1+\kappa_r$ and
$1-\kappa_r$, giving
\begin{equation}
 \det(I_{r-1}+\kappa_rR_{r-1})
   =(1+\kappa_r)^{\lceil(r-1)/2\rceil}
    (1-\kappa_r)^{\lfloor(r-1)/2\rfloor}. \label{eq:reversaldeterminant}
\end{equation}
Both bases are positive because $|\kappa_r|<1$.

\textbf{Compose all recursion steps.} The order-one step is $h^{(1)}=\kappa_1$
and has determinant one. Each later step transforms the earlier coefficients
and appends the next reflection coefficient. Repeated application of the chain
rule multiplies their determinants. Combining with \eqref{eq:jacchain} yields
\begin{equation}
 \mathcal J_m(z_c)=
 \prod_{i=1}^m(1-\kappa_i^2)
 \prod_{r=2}^m
 (1+\kappa_r)^{\lceil(r-1)/2\rceil}
 (1-\kappa_r)^{\lfloor(r-1)/2\rfloor}.
 \label{eq:jacproduct}
\end{equation}

\page
\subsection{Combine the blocks and form the actual sampling target}

Taking logarithms of \eqref{eq:jacproduct} turns every product into a sum:
\begin{equation}
 \begin{aligned}
 \log\mathcal J_m(z_c)
  ={}&\sum_{i=1}^m\log(1-\kappa_i^2)\\
    &+\sum_{r=2}^m\left[
       \left\lceil\frac{r-1}{2}\right\rceil\log(1+\kappa_r)
       +\left\lfloor\frac{r-1}{2}\right\rfloor\log(1-\kappa_r)\right].
 \end{aligned} \label{eq:genericjacobian}
\end{equation}
For $x=(a,b,f,d,s)$ and $z=(z_a,b,f,z_d,v)$, the full derivative is block
diagonal: the $a$ block depends only on $z_a$, the $d$ block only on $z_d$,
and $s$ only on $v$. The $b$ and $f$ blocks are identities and have determinant
one. Equation~\eqref{eq:variancemap} supplies the remaining factor $s=e^v$.
Therefore
\begin{equation}
 \boxed{\log J(z)=\log\mathcal J_{n_a}(z_a)+\log\mathcal J_{n_d}(z_d)+v.}
 \label{eq:jacobian}
\end{equation}
An empty denominator block contributes zero to $\log J$: set $\mathcal J_0=1$.
Unchanged real initial-state coordinates also contribute determinant one.

Let $\ell(\theta,s)$ be the physical log posterior, including the likelihood and
the priors. The numerical target is
\[
 \widetilde\ell(z)=\ell(\theta(z),e^v)
      +\log\mathcal J_{n_a}(z_a)+\log\mathcal J_{n_d}(z_d)+v.
\]
For the Gaussian coefficient priors and stationary-noise likelihood used here,
the formula in the next section expands to
\begin{equation}
 \begin{aligned}
 \widetilde\ell(z)={}&\const-\left(\frac N2+\xi_0\right)v
  -\left(\frac{Q(\theta(z))}{2}+\zeta_0\right)e^{-v}\\
 &-\frac12\log|P_{d(z)}|-\frac12R_{\theta(z)}
  +\log\mathcal J_{n_a}(z_a)+\log\mathcal J_{n_d}(z_d).
 \end{aligned} \label{eq:fullytransformedtarget}
\end{equation}
Here $R_\theta$ is the Gaussian-prior quadratic form defined in the next
section. The physical coefficient of $\log s$ is $-(N/2+\xi_0+1)$; the variance
Jacobian $+v$ changes it to $-(N/2+\xi_0)$ above. Add an initial-state log prior
if those states are unknown. The determinant of $P_d$ belongs to the likelihood;
it does not replace any factor in $J$.

\subsection{What to compute, and when software does the correction}
At each proposed $z$: compute both reflection vectors, run both recursions,
set $s=e^v$, evaluate the physical likelihood and priors, evaluate the two
log Jacobians, and add their sum and $v$ to the physical log posterior.
Differentiate that complete scalar target for HMC.

Automatic differentiation computes derivatives of the expressions supplied;
it does not by itself insert a missing change-of-variables correction.
A package can insert corrections for constraints it transforms internally.
For example, if it handles a declared positive $s$ through a log transform,
its $+v$ is already included; do not add it twice. A user-defined recursion
from $z_a,z_d$ to $a,d$ still needs its own correction [2]. These formulas assume
the explicit $\tanh$ and $\exp$ maps above; use the actual map if it differs.

Assigning Gaussian priors directly to $z_a,z_d$ is a different prior model.
To preserve the physical Gaussian priors, evaluate those priors at $a(z_a)$
and $d(z_d)$ and retain the derived Jacobians. Finally, use stable log expressions
near the boundaries and avoid clipping, which destroys the stated bijection.

\page
\section{The posterior target and one complete evaluation}
\label{sec:target}

Bayes' rule gives
\[
 p(\theta,s\mid\data)\propto L(\theta,s)\,p(a)p(b)p(f)p(d)p(s).
\]
No evaluation of the evidence $p(\data)$ is needed. With zero-mean Gaussian
priors, define
\[
 R_\theta=a^\top V_a^{-1}a+b^\top V_b^{-1}b
                   +f^\top V_f^{-1}f+d^\top V_d^{-1}d.
\]
For the stationary AR noise model of arbitrary order, the log density with
respect to physical coordinates $(\theta,s)$, up to parameter-independent
constants, is
\begin{equation}
 \boxed{\begin{aligned}
 \ell(\theta,s)
   ={}&-\left(\frac N2+\xi_0+1\right)\log s
       -\frac{Q(\theta)/2+\zeta_0}{s}\\
     &-\frac12\log|P_d|-\frac12R_\theta .
 \end{aligned}} \label{eq:physicaltarget}
\end{equation}
Outside the stable region or for $s\leq0$, the density is zero. Fixed prior
normalizing constants can be omitted. If $x_0$ is unknown, use
$Q(\theta,x_0)$ and add $\log p(x_0)$; do not silently fix it when sampling.
For $n_d=0$, interpret $\log|P_d|$ as zero. For a conditional likelihood,
use $Q_{\mathrm{cond}}$ and $N-n_0$, and omit the initial-block determinant.

The numerical sampler targets the density in $z$ coordinates:
\begin{equation}
 \widetilde\ell(z)=\ell\bigl(\theta(z),s(z)\bigr)+\log J(z).
 \label{eq:transformedtarget}
\end{equation}
The stationary-noise determinant term in \eqref{eq:physicaltarget} and the
change-of-variables term in \eqref{eq:transformedtarget} have different origins.
Both are required for this model and this choice of sampling coordinates.

\subsection{What the computer does at each proposed parameter vector}
\begin{enumerate}
 \item Transform $z$ to $(a,b,f,d,s)$ using \eqref{eq:transform}.
 \item Simulate $\alpha_t$ over the fixed input record, using the specified
 initial state and any available prehistory.
 \item Form $\beta_t=\alpha_t+\sum_{j=2}^{n_f}f_j\alpha_t^j$ and
 $e_t=c_t-\beta_t$.
 \item Whiten the residuals, include the chosen initial-noise contribution,
 and compute $Q$ and the log likelihood.
 \item Add coefficient and variance log priors, any initial-state prior,
 and the transformation log Jacobian.
 \item Return $\widetilde\ell(z)$ and its gradient
 $\nabla_z\widetilde\ell(z)$ to the sampler.
\end{enumerate}
Automatic differentiation can propagate through the deterministic recursion.
No separate sampling of the entire $\alpha$, $\beta$ or innovation sequence is
needed for this model. HMC requires first derivatives, not an inverse Hessian.

For a physical-model baseline, the exact IIR recursion is the clearest forward
model. If an FSM approximation is used, its frequency grid, truncation and
boundary convention become part of the likelihood approximation and must be
controlled. The Bayesian calculation should not silently target a different
forward model from the comparison methods.

\page
\section{How one Hamiltonian Monte Carlo transition works}

MCMC constructs a correlated sequence whose stationary distribution is the
posterior. Under appropriate convergence conditions, averages along that
sequence approximate posterior expectations. It does not seek a sequence of
ever smaller losses. A draw with a lower posterior density can be an entirely
appropriate part of the sample.

For a concrete explanation, consider fixed-length HMC [3]. Let $z_k$ be the
current unconstrained position and let $\mathsf M$ be a positive definite
momentum covariance, fixed during retained sampling. Introduce an auxiliary
momentum $p$ and define
\begin{equation}
 p\sim\N(0,\mathsf M),\qquad
 H(z,p)=-\widetilde\ell(z)+\frac12p^\top\mathsf M^{-1}p.
 \label{eq:hamiltonian}
\end{equation}
The momentum is an algorithmic variable, not a state or noise source of the
Wiener system. Its normalization is constant for a fixed $\mathsf M$.

\subsection{Step 1: refresh the momentum}
Draw a new $p_0$ and record $H_0=H(z_k,p_0)$. This random refresh allows the
chain to explore different directions across successive transitions.

\subsection{Step 2: integrate a reversible proposal trajectory}
Starting at $(z_k,p_0)$, take $L$ leapfrog steps of length $\epsilon$. One step is
\begin{align}
 p_{1/2}&=p+\frac\epsilon2\nabla_z\widetilde\ell(z), \notag\\
 z'&=z+\epsilon\mathsf M^{-1}p_{1/2}, \notag\\
 p'&=p_{1/2}+\frac\epsilon2\nabla_z\widetilde\ell(z').
 \label{eq:leapfrog}
\end{align}
Each new gradient evaluation repeats the forward simulation and log-target
calculation from Section~\ref{sec:target}. The leapfrog construction approximately conserves
$H$ and preserves volume. Reverse the final momentum to obtain the standard
reversible proposal; this reversal leaves its kinetic energy unchanged.

\subsection{Step 3: correct numerical integration error}
For the proposal $(z_\star,p_\star)$, accept with probability
\begin{equation}
 a_{\mathrm{acc}}=\min\{1,\exp[H_0-H(z_\star,p_\star)]\}.
 \label{eq:acceptance}
\end{equation}
Generate $U\sim\mathrm{Uniform}(0,1)$. If $U<a_{\mathrm{acc}}$, set
$z_{k+1}=z_\star$; otherwise set $z_{k+1}=z_k$. Discard the auxiliary momentum
and refresh it at the next transition. A rejected transition still contributes
a repeated position to the chain; removing such repetitions changes its law.

\keybox{In exact arithmetic with a correct implementation, the Metropolis
correction accounts for finite leapfrog error. Finite run length and poor
exploration still cause Monte Carlo error. MCMC does not replace the posterior
with a Gaussian approximation, and it does not guarantee that a short run has
explored every posterior mode.}

\page
\section{Use NUTS in practice and run several chains}

Fixed-length HMC illustrates the mechanics, but choosing $L$ by hand is
inconvenient. A mature implementation of the No-U-Turn Sampler (NUTS) is a
reasonable first implementation choice [1]. It builds a trajectory of candidate
points and stops its expansion when continued movement begins to turn back.
The implementation uses a valid candidate-selection rule over the trajectory;
it is not simply fixed-length HMC with its last endpoint accepted by
\eqref{eq:acceptance}.

During warm-up, the implementation adapts the integration step size and a
metric that helps handle different scales and correlations. Retained sampling
then uses the adapted settings. Keep the package's divergence, trajectory
length and acceptance diagnostics. A longer trajectory is more expensive than
a short one, so MCMC iteration counts alone do not measure computation.

\subsection{A concrete initial run protocol}
The following counts are a starting experiment, not a convergence guarantee.
\begin{enumerate}
 \item \textbf{Fix the data and statistical model.} Save the record, input
 history, model orders, initial-state treatment, priors and physical units.
 All chains must target exactly this same posterior.
 \item \textbf{Choose dispersed plausible starting positions.} Use four
 different stable coefficient vectors and positive variances, converted to
 $z$ coordinates. Check that every starting log target and gradient is finite.
 A fitted or MAP estimate can initialize one chain, but all chains should not
 be copies of the same point.
 \item \textbf{Warm up each chain.} For example, use 1,000 warm-up transitions
 per chain. Let the sampler adapt and exclude these transitions from posterior
 summaries.
 \item \textbf{Retain a first batch.} For example, retain 2,000 transitions
 per chain. This yields 8,000 raw draws across four chains, not 8,000 independent
 draws. Preserve chain identities for diagnostics.
 \item \textbf{Assess exploration and precision.} Use the checks in Section~\ref{sec:diagnostics}.
 If exploration is sound but Monte Carlo errors remain too large, run longer.
 If there are unresolved divergences or modes, investigate the geometry,
 initialization or model before interpreting a longer run as a solution.
 \item \textbf{Save enough to reproduce the result.} Store physical draws,
 sampler diagnostics, seeds, warm-up settings, software versions, elapsed time
 and the exact model and configuration used.
\end{enumerate}

\subsection{Two different meanings of burn-in}
\textbf{Data burn-in} propagates the simulated physical filters before retaining
the observation record. It affects the signal and noise initial distributions.
\textbf{MCMC warm-up} adapts the sampler and moves chains away from their
starting points. It does not modify the data or remove an incorrectly modeled
physical transient. These two operations solve different problems.

No new innovation realizations are generated inside a likelihood evaluation.
Randomness during sampling comes from the MCMC algorithm. New physical
innovations are generated later for posterior predictive simulation.

\page
\section{Decide whether the posterior sample is trustworthy}
\label{sec:diagnostics}

Inspect the physical coefficients, $s$, any initial states, and important
derived quantities. Trace and rank plots should show chains exploring the
same distribution rather than remaining separated. Pair plots can reveal
strong dependence, curved ridges or separate solutions.

{\renewcommand{\arraystretch}{1.22}
\begin{tabular}{@{}p{.22\linewidth}p{.72\linewidth}@{}}
\toprule
\textbf{Diagnostic} & \textbf{Interpretation and first action}\\
\midrule
Rank-normalized $\widehat R$ & Compare within-chain and between-chain behavior.
Aim for $\widehat R<1.01$ for the reported quantities. Larger values call for
investigation of poor mixing or different explored regions.\\
Bulk and tail ESS & Account for dependence between draws. With four chains,
a bulk ESS above 400 is a baseline check, not a universal accuracy target.
Tail ESS matters for interval endpoints and tail probabilities.\\
Divergences & Investigate every post-warm-up divergence before final inference.
They can indicate that numerical trajectories fail in an important part of
the posterior. Longer runs alone need not solve this.\\
Trajectory and energy checks & Repeated maximum-tree-depth events indicate
costly exploration. Low energy-transition efficiency can reveal poor geometry.
Check sampler warnings and parameter-pair plots.\\
\bottomrule
\end{tabular}}

These are practical screening rules, consistent with Stan's diagnostic
guidance [4], not a proof that every possible mode has been discovered.

\subsection{Quantify Monte Carlo error separately from posterior uncertainty}
For a scalar quantity $g(\theta,s)$, the posterior standard deviation describes
uncertainty about that quantity given the record. The Monte Carlo standard
error (MCSE) describes numerical uncertainty in its estimated posterior mean.
For a well-mixed stationary sample, a useful interpretation is
\begin{equation}
 \operatorname{MCSE}(\overline g)\approx
 \frac{\operatorname{SD}_{\mathrm{post}}(g)}{\sqrt{\operatorname{ESS}_{\mathrm{mean}}(g)}}.
 \label{eq:mcse}
\end{equation}
Use the software's appropriate mean-specific ESS or direct MCSE estimate;
bulk ESS is a general diagnostic rather than exactly the same statistic [5].

For example, a desired MCSE below $2\%$ of the posterior standard deviation
suggests a mean-specific ESS of roughly 2,500. This illustrates why the
baseline ESS of 400 need not suffice for precise numerical reporting.
Quantile MCSE requires its own calculation.

\subsection{If diagnostics are unsatisfactory}
Check likelihood signs and initial conditions first. Then inspect parameter
scales and posterior dependence, use realistic priors and improve smooth
parameterizations. A smaller integration step may help some divergences, but
does not cure nonidentifiability or guarantee movement between distant modes.
If multiple modes persist, report the issue and consider methods designed for
multimodal exploration. Do not pool chains stuck in different regions as though
their relative frequencies were correct posterior weights.

\page
\section{Turn draws into parameter and prediction results}

After satisfactory diagnostics, let $(\theta^{(m)},s^{(m)})$, $m=1,\ldots,M$,
be the retained physical draws, combining chains only for summaries. For every
coefficient and the variance, report a posterior mean or median, posterior
standard deviation and a $95\%$ equal-tailed credible interval. Include a
correlation plot when blocks compensate for each other.
\begin{equation}
 \widehat\theta_{\mathrm{mean}}=\frac1M\sum_{m=1}^M\theta^{(m)},\qquad
 I_j=[q_{0.025}(\theta_j),q_{0.975}(\theta_j)].
 \label{eq:summaries}
\end{equation}
Transform each draw before summarizing: in general
$\E[\theta(z)]\neq\theta(\E[z])$, and
$\E[\sqrt{s}]\neq\sqrt{\E[s]}$. If a physical-coordinate MAP is needed,
maximize the physical density \eqref{eq:physicaltarget}; maximizing the
transformed density \eqref{eq:transformedtarget} defines a different MAP.

\subsection{Noise-free output uncertainty}
For each posterior draw, propagate the linear block and static nonlinearity to
obtain $\beta_t^{(m)}$. Pointwise quantiles describe uncertainty about the
noise-free response due to uncertain parameters and, where sampled, initial
states. The posterior mean response is the average of these trajectories:
\[
 \widehat\beta_t=\frac1M\sum_m\beta_t^{(m)}.
\]
It generally differs from evaluating the nonlinear forward model at
$\widehat\theta_{\mathrm{mean}}$.

\subsection{Predict a new noisy observation}
For arbitrary AR order and a known next input, a one-step prediction conditional
on an observed history containing the required lags is
\begin{align}
 e_{t+1-i}^{(m)}&=c_{t+1-i}-\beta_{t+1-i}^{(m)},\quad i=1,\ldots,n_d, &
 \nu_{t+1}^{(m)}&\sim\N(0,s^{(m)}), \notag\\
 c_{t+1}^{(m)}&=\beta_{t+1}^{(m)}-\sum_{i=1}^{n_d}d_i^{(m)}e_{t+1-i}^{(m)}
                                      +\nu_{t+1}^{(m)}. \label{eq:predictive}
\end{align}
These predictive intervals include both parameter uncertainty and new
innovation noise, so they have a different interpretation from intervals for
$\beta$. For several future steps, recursively propagate the AR noise and
generate fresh innovations at each step. For an independent repeated record,
also draw its initial noise state under the specified initialization model.

\subsection{Interpretation for the manuscript's numerical studies}
A credible interval is a posterior probability statement under the chosen
model and prior; it is not automatically a frequentist coverage guarantee.
Assess coverage, bias and RMSE by repeating identification on independently
generated datasets. Within one dataset, MCMC draws measure posterior
uncertainty and numerical integration error, not the sampling distribution
across experiments.

Numerical inference does not intrinsically reduce statistical precision when
it accurately explores the intended posterior. Precision still depends on
record length, excitation, model correctness and prior information. Posterior
means can have finite-sample bias, and a poorly mixed numerical approximation
can introduce additional error.

\page
\section{Validate the implementation before a large study}

Start with a short synthetic record and validate identities that have an
independent answer. Knowing simulation truth is useful for debugging and
later performance assessment; it is not information supplied to the estimator.

\subsection{Likelihood and gradient checks}
For stable coefficient vectors of several AR orders, compare the log density
calculated from the dense covariance $sK(d)$ with \eqref{eq:Qstationary}.
They must agree
to numerical precision for the same residual vector and initial-state model.
Separately check the physical-to-unconstrained Jacobian and compare automatic
gradients with central finite differences at interior parameter points.

Check special cases: $d=0$ reduces to independent output noise, $f_j=0$ for
$j\geq2$ gives a linear signal path, and known zero noise prehistory gives
the corresponding conditional innovation recursion. Include tests for the
actual beginning-of-record convention, where indexing mistakes are common.

\subsection{Two analytic conditional distributions for additional checks}
Although the joint posterior for $(a,b,f,d,s)$ is generally nonconjugate,
some full conditional distributions are simple [6]. These identities also
offer optional alternatives to sampling every parameter by HMC.

With $\theta$ and any signal initial state fixed, the inverse-gamma variance
prior gives
\begin{equation}
 s\mid\theta,\data\sim
 \IG\!\left(\xi_0+\frac N2,\;\zeta_0+\frac{Q(\theta)}2\right).
 \label{eq:conditionalvariance}
\end{equation}
Here $N$ corresponds to the full stationary likelihood. Replace it by the
number of conditional residual contributions if using a conditional likelihood.
The independence of the coefficient priors from $s$ matters to this formula.

Next fix $(a,b,d,s)$ and the signal initial state. Define $Z_{t,j}=\alpha_t^{j+1}$,
$j=1,\ldots,n_f-1$, and $w=c-\alpha$. With $f\sim\N(0,V_f)$,
\begin{align}
 S_f&=\left(V_f^{-1}+\frac1s Z^\top K(d)^{-1}Z\right)^{-1}, \notag\\
 m_f&=S_f\frac1s Z^\top K(d)^{-1}w,
 & f\mid a,b,d,s,\data&\sim\N(m_f,S_f). \label{eq:conditionalf}
\end{align}
Use whitening or solves to implement these expressions. This is a Gaussian
conditional for the polynomial coefficients, not a closed-form joint posterior
for the rational and AR coefficients.

\subsection{Optional Gibbs updates}
One can alternate an exact draw from \eqref{eq:conditionalvariance} with a valid
MCMC update of the remaining parameters, or use \eqref{eq:conditionalf} as
another Gibbs block. The resulting hybrid sampler still needs convergence
diagnostics and can mix poorly when blocks are strongly dependent. Joint NUTS
for the complete $(\theta,s)$ vector remains a straightforward first baseline; analytic
conditional updates are an optional refinement, not a prerequisite.

Finally, inspect posterior predictive behavior, whitened residual
autocorrelation and sensitivity to reasonable prior scales. A single successful
fit is not a finite-sample bias or coverage study.

\page
\section{Implementation order, computational cost and references}

\subsection{A practical order of work}
\begin{enumerate}
 \item Fix the forward model and initial conditions; implement and validate
 the scalar innovation likelihood and physical priors.
 \item Add smooth stability transformations, their Jacobian and differentiable
 recursions; validate log targets and gradients.
 \item Run a small four-chain NUTS pilot. Resolve diagnostic problems and
 quantify Monte Carlo errors before increasing record lengths or repetitions.
 \item Save parameter summaries, credible intervals, predictions, diagnostics
 and computational cost for every dataset. Then add this estimator to the
 initialization and record-length experiments using the same observed records
 and explicitly chosen posterior point estimate.
\end{enumerate}

For fixed orders, a forward likelihood evaluation is linear in record length,
with approximate cost
\[
 O\!\left(N(n_a+n_b+n_f+n_d)\right)
\]
when evaluating the polynomial directly.
Differentiation and repeated trajectory steps increase the constant and total
work. A NUTS transition can require many full-record gradient evaluations.
Record wall time, target/gradient evaluations and effective samples per second.
Compare posterior-mean accuracy against GHAM and other estimators on the same
datasets, while reporting the extra uncertainty information MCMC provides.
Optimizer iterations and NUTS transitions are not comparable units of cost.

\subsection{References and further reading}
The likelihoods, stability map and posterior formulas above specialize standard
Gaussian and MCMC calculations to the manuscript's Wiener-AR model. The
following primary sources document the sampling machinery and diagnostics.

{\small
\begin{enumerate}
 \item M. D. Hoffman and A. Gelman (2014). \emph{The No-U-Turn Sampler:
 Adaptively Setting Path Lengths in Hamiltonian Monte Carlo}. Journal of
 Machine Learning Research 15(47), 1593--1623.
 \href{https://www.jmlr.org/papers/v15/hoffman14a.html}{Paper and full text}.

 \item Stan User's Guide. \emph{Reparameterization and change of variables}.
 \href{https://mc-stan.org/docs/stan-users-guide/reparameterization.html}{Official guide}.
 See also the Reference Manual's
 \href{https://mc-stan.org/docs/reference-manual/transforms.html}{constraint transforms}.
 Transformation conventions must match the chosen implementation.

 \item R. M. Neal (2011). \emph{MCMC Using Hamiltonian Dynamics}. In
 Handbook of Markov Chain Monte Carlo, Chapter 5.
 \href{https://arxiv.org/abs/1206.1901}{Author manuscript}.

 \item Stan. \emph{How to Diagnose and Resolve Convergence Problems}.
 \href{https://mc-stan.org/learn-stan/diagnostics-warnings.html}{Official diagnostic guidance}.
 Practical thresholds are checks, not convergence proofs.

 \item Stan Reference Manual. \emph{Posterior Analysis}.
 \href{https://mc-stan.org/docs/reference-manual/analysis.html}{ESS and Monte Carlo error}.

 \item K. P. Murphy (2007). \emph{Conjugate Bayesian Analysis of the Gaussian
 Distribution}.
 \href{https://www.cs.ubc.ca/~murphyk/Papers/bayesGauss.pdf}{Technical note}.
 Gaussian variance and precision conventions; conditional conjugacy should
 not be confused with joint conjugacy for this nonlinear model.

 \item MathWorks. \emph{Reflection coefficients and prediction polynomials}.
 \href{https://www.mathworks.com/help/signal/ref/rc2poly.html}{Step-up recursion}
 and \href{https://www.mathworks.com/help/signal/ref/poly2rc.html}{stability criterion}.
 The Jacobian in this note follows by differentiating the stated recursion.
\end{enumerate}
Sampling references checked on 29 October 2026; coordinate-transformation
references rechecked on 30 October 2026. The illustrative run
lengths and prior scales in this note are proposed starting choices, not
empirical performance results.
}
\end{document}
